\section{The Secrets}

\subsection{What this edition publishes, and the disclosure it owes
the reader first}

This is the hardest problem in the La Salette dossier, and it is
hard in an unusual way: the difficulty is not only historical but
juridical, and the juridical difficulty touches the act of writing
about it. Before any of the evidence is set out, the reader is owed
the disclosure that follows, in full, and before the analysis of
\S6.

The first thing it contains is a distinction that governs everything
afterwards. The \emph{public message} of 19 September 1846---the
discourse examined at \S3, judged by the Bishop of Grenoble on 19
September 1851, preached at the shrine, quoted by popes, and printed
by the custodians---is not the subject of any Roman restriction
whatever. It is approved in its event, freely discussed, and
expounded in this study without reserve. The \emph{secrets} are a
wholly separate stratum: two things said privately to two children,
written for Pius~IX in 1851, and---in one seer's case
only---published in an expanded form thirty-three years after the
event. Everything in the disclosure below concerns the second and
touches nothing in the first.

\begin{statusdossier}{Disclosure: the Roman acts on the secret, their
unresolved force, and this edition's posture}
\dosfield{The governing act}{Supreme Sacred Congregation of the Holy
Office, decree \term{circa vulgo dictum «Secret de la Salette»}, 21
December 1915, \work{AAS} 7 (1915) 594, quoted entire at \S5.7. Its
operative command: \term{eadem Sacra Congregatio mandat omnibus
fidelibus cuiuscumque regionis ne sub quovis praetextu vel quavis
forma, nempe per libros, opuscula aut articulos sive subscriptos
sive sine nomine, vel alio quovis modo, de memorato subiecto
disserant aut pertractent}---the same Sacred Congregation commands
all the faithful of whatever region not to discuss or treat of the
said subject under any pretext or in any form, whether by books,
pamphlets, or articles, signed or anonymous, or in any other way
whatever.}
\dosfield{Its penalties}{Offending priests are to be deprived of any
dignity they hold and suspended by the local ordinary from hearing
confessions and celebrating Mass; lay offenders are not to be
admitted to the Sacraments until they repent; both are further
subject to the sanctions of Leo~XIII's constitution
\work{Officiorum ac munerum} against publishing on religious matters
without the lawful permission of superiors, and of Urban~VIII's
decree of 13 March 1625 against circulating alleged revelations
without the ordinaries' permission.}
\dosfield{Its saving clause}{\term{Hoc autem decretum devotionem non
vetat erga Beatissimam Virginem sub titulo Reconciliatricis vulgo
«de la Salette» nuncupatam}---this decree does not forbid devotion
to the Blessed Virgin under the title of Reconciler commonly called
\term{de la Salette}.}
\dosfield{The acts that applied it}{Holy Office declaration of 12
April 1916 and Sacred Congregation of the Index decree of 5 June
1916, both condemning by name the two-volume \work{La Leçon de
l'Hôpital Notre-Dame d'Ypres --- Exégèse du Secret de la Salette}
(\work{AAS} 8 (1916) 175, 178--179); and the Holy Office decree
\term{Damnatur opusculum} of 9 May 1923, proscribing and condemning
the 1922 Société Saint-Augustin reprint of Mélanie Calvat's text
\term{mandantes ad quos spectat ut exemplaria damnati opusculi e
manibus fidelium retrahere curent}---ordering those whom it concerns
to see that copies are withdrawn from the hands of the faithful
(\work{AAS} 15 (1923) 287--288). All three are quoted or cited at
their loci in \S5.8.}
\dosfield{The unresolved question}{Whether the 1915 command still
binds is \emph{not settled}. The Index of Prohibited Books ceased to
have the force of ecclesiastical law in 1966; the Church's penal law
was recodified in 1917 and again in 1983; the Congregation of the
Index, on whose discipline the decree's penal apparatus partly
rests, was suppressed in 1917. Against that, no act abrogating,
modifying, interpreting, or confirming the decree was located by
this project's sweep of \work{Acta Apostolicae Sedis} volumes~1
(1909) through~98 (2006), volume~75 excepted, or in the document
index of the Dicastery for the Doctrine of the Faith (\S5.9,
\S5.10). That is a bounded negative---the boundary of a search, not
a statement about the world. This edition therefore does \emph{not}
assert that the decree has lapsed, and does \emph{not} assert that
it binds. It reports the dispute.}
\dosfield{Why this edition publishes anyway}{Because the project's
maintainer has expressly directed that this edition, like its
companion study of Fátima, publish and analyze the secret, on the
judgment that the analysis is the value the document exists to add.
The direction is recorded here as what it is: an editorial decision
of a private study project, taken with the juridic question open and
not resolving it.}
\dosfield{What this edition is not}{Nothing in \S5 or \S6 is an
ecclesiastical interpretation, an authentic interpretation of any
canon or decree, a dispensation, a canonical opinion, or advice to
any other author about what that author may lawfully publish. No
ecclesiastical review or approval of this study has been sought,
granted, or is claimed. A reader who concludes that the 1915 decree
binds him should act on that conclusion; this study does not tell
him otherwise.}
\end{statusdossier}

With that on the record, the rest of \S5 does what it did before:
it establishes the acts, the documents, the custody, and the
permissions, because no analysis of a text is worth anything until
the text's transmission is known. \S6 then publishes and analyzes
the text itself.

\subsection{The secret in the earliest record}

That a secret existed is attested from the first interrogations, and
the interrogations are also the earliest evidence for its shape. In
the depositions printed by the 1854 witness, Mélanie, asked whether
the Lady told her nothing else, answered: ``Yes, Sir, but she has
forbidden us to tell it''; asked what it was about, she answered
with a child's exact logic: ``If I tell you what about, you will
soon find out what it is''; asked when it was given, she placed it
``after speaking of the walnuts and the grapes,'' and added that
before giving it the Lady seemed to speak to Maximin and she heard
nothing; asked whether it was in French, she answered that it was
given her ``in Patois.'' Maximin, asked the same question, placed
his after the sentence about grapes and walnuts, said the Lady then
said something to him \emph{in French}, and reported the form of the
prohibition in the plainest possible words: ``You shall not say
this, nor this, nor this.''\endnote{Ullathorne 1854, pp.~44--46,
printing the interrogation answers of both children. The
language difference---Mélanie's secret in \term{patois}, Maximin's
in French---is the children's own testimony as this witness prints
it, and is recorded because it bears on every later claim about the
texts. This study does not resolve it.}

Four facts are therefore in the record from the beginning: there
were \emph{two} secrets, not one; they were given separately and in
different languages; each child was forbidden to tell; and neither
knew the other's. Every later development has to be read against
that starting point.

\subsection{July 1851: two manuscripts for Pius~IX}

In July 1851 the bishop of Grenoble ended the impasse by directing
the children, under obedience, to write their secrets for the Pope.
The 1852 English witness gives the reasoning he used with them---that
visions and revelations of whatever kind ought to be submitted to
the Roman Pontiff, to whom it belongs to judge in these matters---and
the scene: ``they sat down at different tables, and wrote their
respective letters without the smallest hesitation, and exactly as
if they had been copying what they wrote from some original before
them. They signed and sealed their letters, and the Bishop entrusted
them to the vicar-general of his diocese and another priest to carry
to Rome.'' The mandement's own narrative of the same events is
quoted at \S4.3.\endnote{\work{Pilgrimage} 1852, pp.~28--29. The
detail ``as if they had been copying'' is the witness's own
apologetic observation and is quoted as such, not adopted.}

The delivery is dated in the same witness: ``On the 18th of the same
month these precious missives were placed in the hands of the Holy
Father by the persons we have named. His Holiness immediately read
them in the presence of the messengers\ldots\ he said he must read
them again at his leisure, and then added, `These are scourges for
France, but Germany and Italy, and many other countries, deserve the
same.'\,''\endnote{\work{Pilgrimage} 1852, p.~29. The reported papal
sentence is second-hand testimony transmitted by an apologetic
review three months to a year after the fact, and this study reports
it as such---as a \emph{reported} remark, not as a papal act, not as
an approval, and not as evidence of the secrets' content. Reception
of a document in a Roman audience establishes reception, nothing
more.}

A second and more circumstantial account of the same week comes from
the other side of the dossier, and is quoted here because it is the
only description of the physical act of writing that this project
reached. Bishop Zola, writing in 1880 to a French parish priest,
states that Mélanie wrote her secret for the first time on 3 July
1851, at the convent of the Providence at Corenc, by order of the
Bishop of Grenoble and in the presence of two named
men---\term{M.~Dausse, ingénieur en chef des ponts et chaussées} and
\term{M.~Taxis, chanoine de la cathédrale de Grenoble}---and that
she ``filled three large pages at a single stretch, saying nothing
and asking nothing,'' signed without rereading, folded the sheet,
and addressed the envelope ``\term{A Sa Sainteté Pie~IX, à Rome}.''
He adds that on the next day, 4 July, the secret was copied out
again by Mélanie herself at the bishop's house at Grenoble,
\term{dans le but de bien distinguer deux dates des évènements qui
ne doivent pas arriver à la même époque}---in order to distinguish
clearly two dates of events which were not to occur at the same
period---because, having put only one date the first time, she
feared the Pope would not understand and that there would be
ambiguity. He names the couriers: \term{M.~Gérin}, curé of the
cathedral of Grenoble, and \term{M.~Rousselot}, honorary
vicar-general, who on 18 July handed to Pius~IX the letters of the
Bishop of Grenoble and those of Maximin and Mélanie containing their
secrets.\endnote{Zola letter of 24 May 1880, printed in the 1881
Louvain reprint, pp.~27--28, read there on 2026-07-25 and
\emph{collated against the page images} of the same imprint on the
same date (Appendix~A). Two things in this passage are used later
and are flagged here. Zola's ``trois grandes pages'' is the only
statement of the 1851 manuscript's physical extent that this project
reached, and it comes from a promoter writing twenty-nine years
after the fact who does not claim to have seen the manuscript.
Rousselot is the diocesan \term{rapporteur} of \S4.2, so the courier
and the commission's reporter are the same man; the 1852 witness's
``the vicar-general of his diocese and another priest'' is
consistent with Zola's naming and independent of it. The second
writing on 4 July, made to disambiguate \emph{two dates}, is a datum
about the 1851 text's contents---it contained at least two
dates---and is the only such datum in this study.}

There the documentary trail available to this project stops. This
study did not reach the two 1851 manuscripts, any facsimile or
transcription of them, or any Holy See statement about their present
custody or content. Later literature reports that the manuscripts
were found in a Roman archive at the end of the twentieth century
and subsequently published in facsimile; that report was not
verified at any identified witness by this project, and nothing here
rests upon it.\endnote{Recorded negative, bounded: searches for a
Holy See publication, decree, notification, or archival notice
concerning the 1851 La Salette manuscripts produced nothing on the
Dicastery for the Doctrine of the Faith's document index
(\url{https://www.doctrinafidei.va/en/documenti.html}, read
2026-07-25, zero occurrences of ``Salette'') and nothing in the
literal searches of \work{Acta Apostolicae Sedis} described at
\S5.10. The bound is the search, not the world.}

\subsection{The expanded corpus: a documentary chronology}

Between 1851 and 1879 a second body of text appeared, attributed by
Mélanie Calvat to the same 19 September 1846 and progressively
longer. The best witness this project could reach for its
transmission is, ironically, one of its own promoters: a letter of
Salvatore Luigi Zola, Bishop of Lecce, who had been Mélanie's
spiritual director, printed in the 1881 Louvain reprint of the Lecce
booklet. It is a partisan document: it is used here for the
chronology of custody it asserts and for the admissions its author
makes against his own interest, which \S6 analyzes, and never for
its evaluative claims about authenticity, which remain the
promoter's.

\begin{actstable}
1851, 3 July & Mélanie, by order of the Bishop of Grenoble, at the
Providence at Corenc, before Dausse and Taxis & Writes her secret
for the first time: three large pages at one sitting, signed unread,
addressed to Pius~IX & Zola letter, 24 May 1880 \\
1851, 4 July & Mélanie, at the bishop's house, Grenoble & Copies the
secret out again to distinguish \emph{two dates} of events not to
occur at the same period & Zola letter, 24 May 1880 \\
1851, 18 July & Gérin, curé of the cathedral, and Rousselot,
honorary vicar-general & Deliver to Pius~IX the bishop's letters and
the two children's sealed secrets & \work{Pilgrimage} 1852,
pp.~28--29; Zola letter \\
1860 & A director of Mélanie at Marseille & Obtains a manuscript of
her secret & Zola letter, 24 May 1880 \\
1869 & Bishop Petagna of Castellammare di Stabia & That same
document handed to Zola, then Mélanie's spiritual director, by
Petagna's order & Zola letter, 24 May 1880 \\
1870, 30 Jan. & Mélanie & Delivers the same document to Abbé
Félicien Bliard, with her declaration of authenticity and signature,
but ``with small reticences indicated by dots and by \term{etc.}''
for parts she did not yet judge it right to disclose & Zola letter,
24 May 1880 \\
1870, 24 Feb. & Bliard & Sends a certified copy from Nice to a
consultor of the Index at Rome, and likewise to several
ecclesiastical dignitaries & Zola letter, 24 May 1880 \\
1872 & C.-R. Girard, at Grenoble, holding the text from Bliard &
Publishes the secret in \work{Les secrets de la Salette et leur
importance} & Zola letter, 24 May 1880 \\
1873 & Bliard & Publishes the 1870 document with commentary as
\work{Lettres à un ami sur le secret de la bergère de la Salette}
(Naples), with an approval of the curia of Cardinal Sisto Riario
Sforza dated 30 April 1873 & Zola letter, 24 May 1880 \\
1878, end \slash\ 1879 & Leo~XIII & Grants Mélanie a private
audience and, on the promoters' accounts, receives the secret from
her; she remains five months in a convent at Rome writing the
constitutions of the new order the text says the Virgin asked for &
Zola letter, 24 May 1880; Roubaud 1882, p.~10 \\
1878, 21 Nov. & Mélanie, signing ``Marie de la Croix, Victime de
Jésus, née Mélanie Calvat, Bergère de la Salette'' & Dates her
integral text at Castellammare & 1881 Louvain reprint, text-end
subscription \\
1879, 15 Nov. & Vicar general of Lecce & \term{Nihil obstat:
imprimatur}, given at Lecce & 1881 Louvain reprint, imprimatur
formula \\
1879 & Lecce printing & The booklet printed at Lecce at the request
and expense of a private person & Zola letter, 24 May 1880; 1881
reprint's preface \\
1880, 9 May & Mélanie & Writes that the Lecce booklet is out of
print and gives permission for a reprint ``provided there is the
authorization of a bishop'' & 1881 Louvain reprint, letter \\
1880, 14 Aug. & Cardinal Prospero Caterini, Secretary of the Holy
Office & Letter expressing the Holy See's displeasure and the desire
that the booklet be withdrawn from the hands of the faithful,
``if the thing is possible'' (\S5.6) & Roubaud 1882, pp.~8--9 \\
1881 & Louvain, Imprimerie de J.~Lefever & Reprint of the same text
with a preface, two recent letters of Mélanie, and Zola's letter,
carrying no Belgian \term{imprimatur} of its own & The 1881 imprint
itself \\
1882 & Diocesan press and the Paris and provincial Catholic papers &
Public controversy: a communiqué against the secret in the
\work{Semaine religieuse d'Amiens} of 25 June, reprinted in
\work{Le Monde} and elsewhere, and a printed reply by a French
priest & Roubaud 1882, pp.~3--10 \\
1922 & Société Saint-Augustin, Paris--Rome--Bruges & A further
reprint, condemned in 1923 (\S5.8) & \work{AAS} 15 (1923) 287--288
\\
\end{actstable}

Two features of that chronology are decisive for any judgment about
the expanded text, and both come from its own promoter. First, the
text grew: the 1870 document carried deliberate lacunae which, Zola
says, Mélanie ``filled in'' in the later printing. Second, Zola
states plainly that the 1879 text is \emph{not} what went to Pius~IX
in 1851: ``Mélanie n'a pas envoyé à Sa Sainteté Pie~IX tout le
secret qu'elle a publié dernièrement, mais seulement tout ce que la
sainte Vierge lui inspira sur l'heure d'écrire de cet important
document.''\endnote{Zola letter of 24 May 1880 to a French parish
priest, and the passage on the identity of the Lecce text with the
1869 manuscript ``elle a comblé seulement dans ce dernier ces
lacunes, ces petites réticences,'' both printed in \work{Le secret
complet de la Salette, publié par Mélanie de la Salette, Sœur Marie
de la Croix} (Louvain: Imprimerie de J.~Lefever, 1881), pp.~27--31,
read 2026-07-25 in the Internet Archive scan and, on the same date,
\emph{collated line by line against the page images} of that scan
(References; Appendix~A). The 1881 imprint is public domain; its
searchable text layer is poor, and every French sentence quoted from
this volume anywhere in \S5 and \S6 was checked against the page
image before being printed here.} The 1879 text is
therefore, on the testimony of the bishop who procured its printing,
a different and larger document from the 1851 manuscript. Whatever
else is disputed, that is not.

\subsection{What the 1879 permission was}

The permission under which the Lecce booklet appeared is quoted
here exactly, because a great deal has been built on it.

\begin{operativetext}{Printing permission at the end of the Lecce
text, as reproduced in the 1881 Louvain reprint}
\term{Nihil obstat: imprimatur.}\par
\term{Datum Lycii ex Curia Ep. die 15 Nov. 1879.}\par
\term{Vicarius Generalis, Carmelus Arch.\ Cosma.}\endnote{1881
Louvain reprint, at the end of the reproduced text, immediately
after Mélanie's subscription dated ``Castellamare, le 21 Novembre
1878.'' \term{Lycium} is the Latin name of Lecce; \term{ex Curia
Ep[iscopali]} is the episcopal curia. \textbf{A correction to this
project's earlier record.} An earlier state of this study reported a
material discrepancy between this formula (15 November 1879) and
Zola's letter of 24 May 1880 in the same volume, which the
searchable text layer renders ``a été bien donnée le 18 novembre
1879,'' and preserved the two dates as unresolved. Both pages were
collated against the imprint's page images on 2026-07-25: the
formula at p.~22 reads \emph{15 Nov.\ 1879} and Zola's letter at
p.~31 reads \emph{le 15 novembre 1879}. The two witnesses agree; the
apparent discrepancy was a defect of the text layer and is recorded
here, and in \work{research/source-audit.md}, as a correction rather
than removed. The 1879 Lecce imprint itself was still not reached,
so the formula remains a reprint's reproduction of a permission
(Appendix~A).}
\end{operativetext}

An \term{imprimatur} of this kind is a local censorship permission:
a declaration by a diocesan curia that it found nothing in the text
contrary to faith or morals such as to prevent printing. In the
discipline of the period it was granted by the ordinary of the place
of printing or of the author's residence, at the request of the
person publishing---here, as Zola says, ``the person who wished to
publish it at his own expense and according to his pious
intentions.'' It is not a judgment that the text is authentic, that
its author saw or heard what she reports, or that the Church accepts
its contents. Article~4 of the 1851 mandement had reserved
devotional publications concerning La Salette to the written
approval of the Bishop of \emph{Grenoble}; a permission obtained at
Lecce, thirty-three years later, in another country, from another
curia, is not that approval, and the 1851 act plainly did not
contemplate it.\endnote{The description of \term{imprimatur} in the
period discipline is editorial synthesis, based on the two acts this
study does quote: the 1915 decree's own reference to Leo~XIII's
constitution \work{Officiorum ac munerum} against those who publish
books on religious subjects without the lawful permission of
superiors, and to Urban~VIII's decree of 13 March 1625 against those
who circulate alleged revelations without the ordinaries' permission
(\S5.7). Neither constitution was read in an edition by this project
(Appendix~A).}

Two further features of the permission belong here, because the
promoters' own accounts state them. Zola says in the same letter
that he himself had the text examined by his episcopal curia
``according to the rules of the Church,'' and that his vicar
general, finding no reason to oppose publication, issued the licence
to print---so the reviewing authority was the curia of the author's
own former spiritual director, at the request of a private person
paying for the printing. And the anonymous \term{avant-propos} of
the 1881 reprint asserts something stronger than the document
supports: ``\term{Mgr Zola, évêque de Lecce, lui a donné
l'imprimatur}''---Bishop Zola gave her the imprimatur---where the
printed formula is subscribed by the vicar general and Zola's own
letter attributes the licence to that vicar general. The inflation
is small and it is characteristic.\endnote{Zola letter, 1881 Louvain
reprint, p.~31 (``Je l'ai moi-même fait examiner par ma curie
épiscopale, suivant les règles de l'Eglise, et mon vicaire général,
n'ayant trouvé aucune raison qui pût s'opposer à la publication du
secret, a délivré sa licence d'imprimer en ces termes: «~Nihil
obstat, Imprimatur~», à la personne qui voulait le publier à ses
frais et selon ses pieuses intentions''), and the reprint's
\term{avant-propos}, unpaginated leaf following the title page, n.~2;
both collated against the page images on 2026-07-25. The
\term{avant-propos} is unsigned and is the reprint editor's, not
Mélanie's and not Zola's; this study attributes it to the reprint.
The 1881 volume carries no Belgian \term{imprimatur} of its own: in
place of one, its title-page verso prints Mélanie's own formula,
``Je soumets cette publication au jugement du Saint-Siège
Apostolique et je déclare condamner et rétracter à l'avance tout ce
qu'il y trouverait de contraire à la doctrine Catholique.'' That is
a submission, not a permission, and it is worth noticing against her
own condition of 9 May 1880 that a reprint be made only ``provided
there is the authorization of a bishop.''}

\subsection{The letter of 14 August 1880}

Between the Lecce printing and the Roman decrees of the next century
there is one further intervention, and this study reached it only in
2026, through a witness that reports it while trying to explain it
away. In 1882 a French priest, Abbé I.~Roubaud, published at the
same Louvain press a defence of the secret against a communiqué that
had appeared in the \work{Semaine religieuse d'Amiens} on 25 June
1882 and had been reprinted in \work{Le Monde} and other papers. The
communiqué asserted, he says, ``that the booklet has been condemned
at Rome.'' In answering it, Roubaud concedes the following---and a
defender's concessions are the most useful sentences in a partisan
pamphlet.

\begin{witnesstext}{Abbé I.~Roubaud, \work{La vérité sur le secret
de Mélanie} (Louvain: Lefever, 1882), pp.~8--9}
«~Quelques rares opposants désiraient cette condamnation.
\emph{L'Index s'y refusa.} L'Inquisition fut saisie de l'affaire,
parce qu'il s'agissait d'une \emph{révélation}.~»\par
\medskip
«~Il fallait bien pourtant contenter l'opposition, lui donner un os
à ronger. Voilà pourquoi la lettre du C.~Catérini, secrétaire de
l'Inquisition, vit le jour. Elle est du 14 août 1880.~»\par
\medskip
«~Cette lettre ne renfermait aucun ordre de la livrer à la
publicité; c'était une lettre privée. Pourtant on la publia. Nous
n'en eûmes jamais que la traduction; nous n'en vîmes point le texte
original (le texte latin), ni partiel, ni complet.~»\par
\medskip
«~Cette lettre n'exprime, en effet, que le \emph{déplaisir} du
S.~Siége.~»\par
\medskip
«~Vous me direz peut-être, Monsieur: --- Si le S.~Siége croyait au
Secret publié par Mélanie, le C.~Catérini n'aurait pas exprimé le
désir de voir retirer l'opuscule \emph{des mains des fidèles}. Je
réponds: 1\textsuperscript{o} Son Eminence le C.~Catérini ajouta:
«~si la chose est possible~».~»\endnote{\work{La vérité sur le
secret de Mélanie ex-bergère de La Salette}, par l'abbé I.~Roubaud
(Louvain: Imprimerie de J.~Lefever, 1882), 2nd~edition, pp.~8--9,
read 2026-07-25 in the Internet Archive scan and collated against
its page images the same day (References). The pamphlet is an
openly partisan defence, printed by the same house as the 1881
reprint, carrying no \term{imprimatur} but the author's own
submission formula; it is used here \emph{only} for the facts it
concedes against its own case and for the state of the 1882
controversy, never for its arguments or its evaluations. Prospero
Caterini (1795--1881) was Secretary of the Supreme Congregation of
the Holy Office. Roubaud's parenthesis that Caterini ``was only the
secretary'' because the Prefect of the Holy Office is the Pope is
correct as to the office. This project did not reach the letter
itself in any form, Latin or vernacular, and does not quote it: what
is reported here is that a defender of the secret, writing two years
after the fact and printing his address, states that such a letter
exists, gives its date, states its author and office, states its
substance, and states that he had never seen its original.}
\end{witnesstext}

What that passage establishes, at the ceiling of a partisan witness
reporting against his own interest, is a short list and nothing
more: that in 1880 some sought a condemnation of the booklet at
Rome; that the Congregation of the Index declined to give one; that
the Holy Office took the matter up because a \emph{revelation} was
involved; that a letter went out over the signature of the Secretary
of the Holy Office on 14 August 1880; that it expressed the Holy
See's displeasure and the desire that the booklet be withdrawn from
the hands of the faithful, qualified by ``if the thing is
possible''; and that it was a private letter which the defenders
knew only in a translation whose Latin original they had never seen.

Three cautions attach, and this study attaches all three. A private
letter of a congregation's secretary is not a decree, is not an act
of the Apostolic See published in its gazette, and is not law; the
literal search of \work{Acta Sanctae Sedis} for 1879--1881 described
at \S5.10 found no such act and no occurrence of Caterini's name in
those volumes. The letter's wording is not established: no text of
it, in any language, was reached by this project, and Roubaud's
quotation of the qualifying phrase is a defender's quotation of a
translation. And a request to withdraw a book ``if possible'' is a
prudential measure, not a censure of a proposition. What it is
\emph{not} is nothing: the earliest reachable trace of the Holy
See's own mind about the Lecce booklet, ten months after it
appeared, is displeasure and a wish to see it withdrawn.

\subsection{The decree of 21 December 1915}

Thirty-six years after the Lecce printing, the Supreme Sacred
Congregation of the Holy Office issued the act that still governs
the subject. It was published in the official gazette, and this
study read it there.

\begin{operativetext}{Suprema Sacra Congregatio S. Officii, decree
\term{circa vulgo dictum «Secret de la Salette»}, 21 December 1915:
\work{Acta Apostolicae Sedis} 7 (1915) 594}
Ad Supremae huius Congregationis notitiam pervenit quosdam non
deesse, etiam ex ecclesiastico coetu, qui, posthabitis responsionibus
ac decisionibus ipsius S.~Congregationis, per libros, opuscula atque
articulos in foliis periodicis editos, sive subscriptos sive sine
nomine, de sic dicto \term{Secret de la Salette}, de diversis ipsius
formis, nec non de eius praesentibus aut futuris temporibus
accommodatione disserere ac pertractare pergunt; idque non modo
absque Ordinariorum licentia, verum etiam contra ipsorum vetitum. Ut
hi abusus qui verae pietati officiunt, et ecclesiasticam auctoritatem
magnopere laedunt, cohibeantur, eadem Sacra Congregatio mandat
omnibus fidelibus cuiuscumque regionis ne sub quovis praetextu vel
quavis forma, nempe per libros, opuscula aut articulos sive
subscriptos sive sine nomine, vel alio quovis modo, de memorato
subiecto disserant aut pertractent.\par
Quicumque vero hoc Sancti Officii praeceptum violaverint, si sint
sacerdotes, priventur omni, quam forte habeant, dignitate et per
Ordinarium loci ab audiendis sacramentalibus confessionibus et a
missa celebranda suspendantur: et si sint laici ad Sacramenta non
admittantur donec resipiscant. Utrique insuper subiaceant
sanctionibus latis tum a Leone PP.~XIII per Constitutionem
\work{Officiorum ac munerum} contra eos qui libros de rebus
religiosis tractantes sine legitima Superiorum licentia publicant,
[tum] ab Urbano~VIII per decretum \work{Sanctissimus Dominus Noster}
datum die 13 martii 1625 contra eos qui assertas revelationes sine
Ordinariorum licentia vulgant.\par
Hoc autem decretum devotionem non vetat erga Beatissimam Virginem
sub titulo \term{Reconciliatricis} vulgo \term{de la Salette}
nuncupatam.\par
Datum Romae, ex Aedibus Sancti Officii, die 21 decembris
1915.\endnote{\work{Acta Apostolicae Sedis} 7 (1915), p.~594, read
2026-07-25 in the Holy See's archive PDF
(\url{https://www.vatican.va/archive/aas/documents/AAS-07-1915-ocr.pdf}),
which is a typeset text transcription rather than a page facsimile;
it carries no page images, so no image collation was possible
(Appendix~A). Two evident defects of that transcription are
corrected in square brackets or recorded: the second correlative of
\term{tum\ldots tum} is printed ``cum,'' bracketed here as [tum];
and the saving clause is printed ``Beconciliatricis,'' read here as
\term{Reconciliatricis}, the reading required by the title's
universal form and confirmed by the same title in the 1961 papal
letter quoted at \S9.4. The volume's own analytical index confirms
the locus, entering the act at ``Dec.~21 Decretum circa vulgo dictum
«Secret de la Salette», 594'' and, under \term{Salette}, ``prohibetur
a S.~C.\ldots\ 594.'' Subscribed \term{Aloisius Castellano, S.~R.\ et
U.~I.\ Notarius}.}
\end{operativetext}

In editorial rendering: it has come to the knowledge of this Supreme
Congregation that there are some, even from the ecclesiastical body,
who---disregarding the responses and decisions of this same Sacred
Congregation---continue, through books, pamphlets, and articles
published in periodicals, whether signed or anonymous, to discuss
and treat of the so-called \term{Secret de la Salette}, of its
various forms, and of its application to present or future times,
and this not only without the permission of the ordinaries but even
against their prohibition. That these abuses, which harm true piety
and gravely injure ecclesiastical authority, may be restrained, the
same Sacred Congregation commands all the faithful of whatever
region not to discuss or treat of the said subject under any pretext
or in any form---by books, pamphlets, or articles, signed or
anonymous, or in any other way whatever. Priests who violate this
precept of the Holy Office are to be deprived of any dignity they
may hold and suspended by the local ordinary from hearing
sacramental confessions and from celebrating Mass; lay persons are
not to be admitted to the Sacraments until they repent; and both are
further subject to the sanctions laid down by Leo~XIII in
\work{Officiorum ac munerum} against those who publish books
treating of religious matters without the lawful permission of
superiors, and by Urban~VIII in the decree of 13 March 1625 against
those who circulate alleged revelations without the permission of
the ordinaries. This decree does not forbid devotion to the Blessed
Virgin under the title of Reconciler commonly called \term{de la
Salette}.

Four things in that text are worth marking. It is aimed at
\emph{discussion and treatment in publication}, and its stated
grievance is disobedience to ordinaries. It binds ``all the faithful
of whatever region''---a universal precept, not a local one. It
carries penalties that are, in form, penal-law sanctions of a legal
system that has since been replaced twice. And it ends by expressly
protecting the devotion: whatever the Holy Office thought of the
secret literature, it did not touch Our Lady of La Salette.

\subsection{The acts of 1916 and 1923}

The 1915 decree did not stand alone. Two further acts, both verified
in the official gazette, show the Holy See applying it to named
publications.

\textbf{1916.} On 12 April 1916 the Holy Office declared a
two-volume work condemned and proscribed under the general rules of
\work{Officiorum ac munerum}: \work{La Leçon de l'Hôpital
Notre-Dame d'Ypres --- Exégèse du Secret de la Salette}, by
``Dr~Henry Mariavé,'' volume~I, Paris 1915; volume~II,
\work{Appendices}, Montpellier 1915. The declaration is dated 13
April 1916. On 5 June 1916 the Sacred Congregation of the Index
condemned and proscribed the same work by name in its own decree,
which Benedict~XV approved and ordered promulgated on 6
June.\endnote{\work{AAS} 8 (1916), p.~175 (Holy Office,
\term{Declaratio circa opus quoddam}, Feria~IV, 12 April 1916, dated
13 April 1916) and pp.~178--179 (Sacred Congregation of the Index,
decree of 5 June 1916, listing the work with the note ``Decr.\
S.~Off.\ 12 apr.\ 1916''; the index decree is signed
Fr.~Card.~Della Volpe, Prefect, and Thomas Esser,~O.P., Secretary,
and dated 6 June 1916), read 2026-07-25 in the Holy See's archive
PDF
(\url{https://www.vatican.va/archive/aas/documents/AAS-08-1916-ocr.pdf}).
The author's name appears as ``Henry Mariavé'' in the Holy Office
declaration and ``Henri Mariavé'' in the Index decree; both spellings
are as printed. The work itself was neither sought nor consulted,
and nothing about its contents is asserted.}

\textbf{1923.} On 9 May 1923 the Holy Office proscribed and
condemned the 1922 reprint of Mélanie's text and ordered its
withdrawal from the hands of the faithful.

\begin{operativetext}{Suprema Sacra Congregatio S. Officii, decree
\term{Damnatur opusculum}, Feria~IV, 9 May 1923: \work{Acta
Apostolicae Sedis} 15 (1923) 287--288}
In generali consessu Supremae Sacrae Congregationis S.~Officii Emi
ac Rmi Domini Cardinales fidei et moribus tutandis praepositi
proscripserunt atque damnaverunt opusculum: \work{L'apparition de la
très Sainte Vierge sur la sainte montagne de la Salette le samedi 19
septembre 184[6]. --- Simple réimpression du texte intégral publié
par Melanie, etc.} Société Saint-Augustin, Paris--Rome--Bruges,
1922; \term{mandantes ad quos spectat ut exemplaria damnati opusculi
e manibus fidelium retrahere curent.}\par
Et eadem feria ac die Sanctissimus D.~N.~Pius divina providentia
Papa~XI, in solita audientia R.~P.~D.~Assessori S.~Officii impertita,
relatam sibi Emorum Patrum resolutionem approbavit.\par
Datum Romae, ex aedibus S.~Officii, die 10 maii 1923.\endnote{\work{AAS}
15 (1923), pp.~287--288, read 2026-07-25 in the Holy See's archive
PDF
(\url{https://www.vatican.va/archive/aas/documents/AAS-15-1923-ocr.pdf}),
again a typeset transcription without page images. The archive text
prints the pamphlet's title-year as ``1845''; since the apparition
fell on Saturday 19 September 1846 and 19 September 1845 was not a
Saturday, the year is given here as 184[6] with the transcription's
reading recorded rather than silently replaced, and no page image
was reachable for collation (Appendix~A). The volume's index
confirms the locus at ``Salette, 287'' and ``Salettensis apparitio,
288.'' Subscribed \term{Aloisius Castellano, Supremae S.~C.\
S.~Officii Notarius}. The decree immediately preceding it in the same
gazette section concerns unrelated works and is not part of this
dossier.}
\end{operativetext}

In editorial rendering: in the general session of the Supreme Sacred
Congregation of the Holy Office, the Cardinals charged with guarding
faith and morals proscribed and condemned the pamphlet
\work{L'apparition de la très Sainte Vierge sur la sainte montagne
de la Salette le samedi 19 septembre 1846 --- Simple réimpression du
texte intégral publié par Mélanie}, published by the Société
Saint-Augustin at Paris--Rome--Bruges in 1922, ordering those whom
it concerns to see that copies of the condemned pamphlet are
withdrawn from the hands of the faithful; and on the same day
Pius~XI approved the Cardinals' resolution.

Note precisely what the 1923 act condemns: a named 1922 \emph{reprint
of the text published by Mélanie}. It is a book condemnation, of a
specific imprint, with a recall order. It does not mention the
apparition, the 1851 judgment, the shrine, the devotion, or the
Missionaries, and it does not purport to judge the 1851 manuscripts.

\subsection{The standing juridic question}

The question a serious reader will ask is whether the 1915 command
still binds. This study does not answer it, and here is why the
answer is genuinely unsettled rather than merely unknown to the
author.

\begin{boundarytable}{What the acts settle}{What remains unsettled}
The 1915 decree issued a universal precept of the Holy Office to all
the faithful, forbidding publication and treatment of the subject.
Its own text presents this as a command (\term{mandat}), distinct
from the penalties that follow it. & Whether a non-penal precept of
the Supreme Congregation, issued under the discipline of 1915,
survives the two subsequent recodifications of the Church's law is a
question of canonical interpretation that this study is not
competent to settle and that no act located here settles. \\
The penal apparatus attached to it is expressly built on the
sanctions of \work{Officiorum ac munerum} and on the Index
discipline. & Those supports have gone. The Sacred Congregation of
the Index was suppressed in 1917; the Index of Prohibited Books
ceased to have the force of ecclesiastical law by the notification
of 14 June 1966, published in \work{AAS} 58 (1966) 445; and the 1983
Code retains no penalty in these terms. \\
The 1923 act condemned a specific 1922 imprint and ordered its
recall. & A condemnation of an imprint does not survive as a
prohibition of a book once the Index discipline it belonged to has
lost legal force; equally, the lapse of that discipline does not
convert a condemned text into an approved revelation. \\
The 1915 decree expressly preserves devotion to Our Lady under the
title \term{Reconciliatrix} of La Salette. & Nothing in the 1915,
1916, or 1923 acts touches the 1851 judgment, the cult, the shrine,
or the Missionaries, all of which have continued under later Roman
acts (\S8, \S9). \\
\end{boundarytable}

What can be said without controversy is narrower and is enough for
practice: no act located by this project has ever declared any
version of the secret to be of supernatural origin; the Holy See's
public interventions on the subject were, without exception,
restrictive; the devotion was expressly and repeatedly protected;
and a Catholic who simply sets the secret literature aside is in no
difficulty whatever. Two further things follow, and this edition
says both against itself. Anyone who circulates the 1879 text as a
revelation to be believed, a chronology to be applied, or a warrant
for judging the Church's pastors is acting against the manifest mind
of every act that exists. And a study which publishes and analyzes
that text---as \S6 does, on the disclosure at \S5.1---is doing
something the 1915 decree's words, read at their widest, would also
cover; that is precisely why the disclosure stands where it stands,
and why this study neither claims the decree has lapsed nor pretends
it is not there.\endnote{The 1966
notification is \work{AAS} 58 (1966), p.~445, whose locus this study
takes from the repository's existing source record for that act and
did not itself re-verify in the gazette; every other AAS and ASS
locus in this study was read directly in the Holy See's archive PDFs
on 2026-07-25 (Appendix~A). The suppression of the Congregation of
the Index in 1917 is stated here as common institutional history and
no source is claimed for it. The final sentence is editorial
judgment, labelled as such, and is not a canonical opinion.}

\subsection{Bounded negatives}

Because the temptation to fill this record by inference is strong,
the searches that came back empty are recorded
exactly.\endnote{Search record: the AAS sweep, its exact volume
range, its method, and its per-volume hit counts are recorded in
\work{research/source-audit.md}; the ASS volumes searched are 11
(1878), 12 (1879), 13 (1880), and 14 (1881). All searches were run
2026-07-25. A literal search over an OCR-derived text layer
establishes that the string was not found in those bytes; it does
not establish that no act exists.}

\begin{itemize}[leftmargin=1.2em,itemsep=0.15em,topsep=0.25em]
\item \textbf{Later Roman acts.} The annual volumes of \work{Acta
Apostolicae Sedis} from volume~1 (1909) through volume~98 (2006)
were downloaded from the Holy See's archive and searched literally
and case-insensitively for the string \emph{salette} in the
archive's text layer---every volume in that range except volume~75
(1983), which that route does not serve; volumes from 99 (2007)
onward are not published there as single annual files and were not
searched. After 1923 the hits are of three kinds only: routine
personnel and jurisdiction notices concerning the Missionaries of
Our Lady of La Salette; the approval of that congregation's
constitutions (\S8.6); and the two substantive Marian items quoted
at \S9.3 and \S9.4. \emph{No act modifying, interpreting,
confirming, or abrogating the 1915 decree was found.} The search is
literal and cannot exclude OCR damage, paraphrase, a differently
spelled heading, an unsearched volume, or an act published elsewhere
than in the gazette.
\item \textbf{The Dicastery's index.} The document index of the
Dicastery for the Doctrine of the Faith contains no La Salette entry
(read 2026-07-25; zero occurrences of the string).
\item \textbf{The 1880 letter.} The letter of 14 August 1880 reported
at \S5.6 was not reached in any form: not in \work{Acta Sanctae
Sedis} for 1878--1881, whose volumes contain no occurrence of
Caterini's name at all; not in any Latin or vernacular printing; and
not at any archival notice. Its existence, date, author, and
substance rest wholly on one partisan witness of 1882 who states
that he himself saw only a translation. Nothing in \S6 rests on its
wording.
\item \textbf{The 1879 Lecce imprint.} No copy, scan, or page image
of the 1879 Lecce printing was reached. The permission formula, the
subscription, and the whole of the text analyzed in \S6 are read in
the 1881 Louvain reprint, collated against that imprint's page
images; the relation of the 1881 reprint to the 1879 Lecce setting
rests on Zola's statement and on the reprint's own claim, and was
not itself verified (\S6.2).
\item \textbf{The 1851 manuscripts.} Not reached in any form
(\S5.3). This is the decisive gap in the whole dossier and \S6.9
states what its closure would do to the judgment reached there.
\item \textbf{Maximin's secret.} No printing, transcription, or
facsimile of any text purporting to be Maximin Giraud's secret was
reached at an identified imprint. Texts circulate; none was
identified to an edition this project could name and check, and
none is published, quoted, or characterized here (\S6.10).
\end{itemize}
